\documentclass[11pt]{article}
\usepackage[T1]{fontenc}
\usepackage{lmodern,amsmath,amssymb,amsthm,mathtools,microtype,booktabs}
\usepackage[margin=1.05in]{geometry}
\usepackage{xcolor}
\usepackage[colorlinks=true,linkcolor=blue!50!black,citecolor=blue!50!black,urlcolor=blue!50!black]{hyperref}
\newtheorem{theorem}{Theorem}[section]
\newtheorem{lemma}[theorem]{Lemma}
\newtheorem{corollary}[theorem]{Corollary}
\newtheorem{proposition}[theorem]{Proposition}
\theoremstyle{remark}\newtheorem{remark}[theorem]{Remark}
\newcommand{\E}{\mathbb E}
\newcommand{\Pp}{\mathbb P}
\newcommand{\G}{\mathcal G}
\newcommand{\Pool}{\mathcal P}
\newcommand{\DGRH}{\textup{(DGRH)}}
\newcommand{\abs}[1]{\left|#1\right|}
\newcommand{\floor}[1]{\left\lfloor #1\right\rfloor}
\newcommand{\ceil}[1]{\left\lceil #1\right\rceil}
\newcommand{\Rea}{\operatorname{Re}}
\newcommand{\ind}{\mathbf 1}
\newcommand{\normcount}{\frac{\varphi(q)R_{S,k}(g)}{\binom mk}}
\title{The quadratic-fibre transition for squarefree products of small primes\\[3pt]\large A sharp finite-sampling law at the logarithm-squared cutoff}
\author{Research note for the ProveIt project}
\date{27 November 2026}
\setlength{\emergencystretch}{3em}
\begin{document}
\maketitle
\begin{abstract}
Assume the Dirichlet Riemann hypothesis exported by the vendored \texttt{grh-lean} development. For any fixed $a>1/3$ and sufficiently large $q$, let $S$ be an arbitrary retained fraction at least $a$ of the primes $p\le C\log^2 q$ coprime to $q$. Put $m=|S|$, choose exactly $k$ distinct primes, and write $h=\min(k,m-k)\ge B\log q$. Let $t$ be the minimum number of retained primes outside a single fibre of a nontrivial quadratic character. We show that, after a phase correction, the normalized counts of products are uniformly equal to $1\pm (1-2h/m)^t$, with a relative error $O(1/h+q^{-A})$ on both occupied sign classes for any fixed $A>0$. The missing class when $t=0$ is an exact obstruction. In particular, at sparse quotas $h=o(m)$, equidistribution holds if and only if $ht/m\to\infty$. The sparse-quota transition law is sharp: when $ht/m\to\lambda>0$, the bias tends to $e^{-2\lambda}$, and the first finite-population correction to the binomial bias is $-2\lambda^2e^{-2\lambda}/h+o(1/h)$. The proof isolates the only potentially large Fourier mode and establishes a uniform, exponentially localized parity approximation for hypergeometric sampling. This answers the near-one-half density question in the merged ProveIt article in a quantitative asymptotic sense, including rare-fibre relative counts.
\end{abstract}

\section{The input and the question}
The input is the precise zero assertion
\begin{equation}\tag{DGRH}\label{eq:dgrh}
 L(s,\chi)\ne0\quad(\Rea s>1/2),
 \qquad\text{except for the principal-character pole at }s=1,
\end{equation}
for all Dirichlet characters of all positive moduli, together with the corresponding critical-strip zero statement. The vendored readable edition of \texttt{grh-lean} \cite{Lean} is identified by its README as a byte-exact import of upstream revision
\begin{center}\small\texttt{72fead102f663d5c7a407964a13dc0cf9d977207}.\end{center}
We inspected the source files and the article's TeX (not its compiled PDF). Specifically, \texttt{src/Statements.lean} defines \path{DirichletLNonvanishingRightOfHalf}\newline and \texttt{DirichletLZerosOnLine} using Mathlib's functions, and \texttt{src/Main.lean} exports \path{LFunction_ne_zero_of_half_lt_re} and \texttt{LFunction\_eq\_zero\_re}. The former is exactly \texttt{New.grh}; the latter follows using reflection of zeros in \texttt{New/Main.lean}. The main arithmetic inequality in \path{New/Reduction/MainInequality.lean} compares corrected covolumes; \path{New/Reduction/OneSidedBound.lean} turns it into the half-plane statement through a smoothed prime remainder and a zero-exclusion criterion. We use the exported conclusions, not a uniform error estimate from the internal admissible-scale argument. The uniform character bounds below follow independently from classical explicit-formula estimates once \eqref{eq:dgrh} is available.

The upstream README reports the standard three Lean axioms, \texttt{propext}, \texttt{Classical.choice}, \texttt{Quot.sound}. Sutherland's independent verification report for the same revision records successful sandboxed builds, proof checks and statement comparison \cite{Suth}. The present work is a downstream mathematical proof, not an independent replay of the entire Lean project. No zero simplicity or linear independence hypothesis is needed; no non-Dirichlet GRH is invoked. The new argument is not itself kernel-checked. We label the arithmetic conclusions as consequences of DGRH to record their dependency, rather than asserting a new independent proof of the input.

The merged article \cite{Article} assembles 131 reports. Its section \texttt{dir.tex}, under ``Squarefree products of small primes after deletions'', establishes a bounded quotient obstruction (label \texttt{dir:thm:quotient}), a fixed-cardinality coefficient estimate (\texttt{dir:lem:circle}), and uniform small bias at $C\log^2 q$ (\texttt{dir:lem:bias}). Its open-questions subsection asks whether distance from a quadratic fibre can replace a fixed retention margin above $1/2$. We refine this exact quotient law by proving a universal first-order approximation and an optimal transition scale, even below density $1/2$. The article already gives the exact quadratic coefficient as a signed binomial sum: that identity and support alone are not claimed as new here.

\section{Statement of the result}
Let $q\ge3$, $L=\log q$, $G=(\mathbb Z/q\mathbb Z)^\times$, and
\[
 \Pool=\{p\le y:p\text{ prime},\ p\nmid q\},\quad y=C L^2,\quad N=|\Pool|.
\]
For $S\subseteq\Pool$, put $m=|S|$ and
\[
 R_{S,k}(g)=\#\bigl\{T\subseteq S:|T|=k,\ \prod_{p\in T}p\equiv g\pmod q\bigr\}.
\]
All primes are distinct labels, even if they have equal residues. A quadratic character here means a nontrivial homomorphism $G\to\{\pm1\}$. Such a character exists because $q\ge3$ implies $G$ has even order. Define
\begin{equation}\label{eq:distance}
 t=\min_{\chi^2=1,\ \chi\ne1}\min_{\sigma\in\{\pm1\}}
       \#\{p\in S:\chi(p)\ne\sigma\}.
\end{equation}
Choose any minimizing pair $(\chi_*,\sigma)$ (ties are allowed). Set
\[
 h=\min(k,m-k),\quad u=h/m,\quad \lambda=ut,\quad b=1-2u,
 \quad s_k=\sigma^k(-1)^{t\ind_{k>m/2}}.
\]
We use $0^0=1$ in $b^t$. The phase $s_k$ matters for complementary quotas.

\begin{theorem}[Uniform quadratic-fibre law; uses DGRH]\label{thm:main}
Fix $a>1/3$ with $a\le1$ and $A>0$. There exist $B,C>0$ depending only on $a,A$ such that, for all sufficiently large $q$, all $S\subseteq\Pool$ with $m\ge aN$, and all $k$ with $h\ge B L$, the following hold.
\begin{enumerate}
\item If $t=0$, no product belongs to the sign class $s_k\chi_*(g)=-1$. For every $g$ in the other class,
\[
 \normcount=2+O(q^{-A}).
\]
\item If $t\ge1$, every $g\in G$ occurs, and uniformly over $q,S,k,g$,
\begin{equation}\label{eq:relative}
 \normcount=\bigl(1+s_k\chi_*(g)b^t\bigr)
     \left(1+O\left(\frac1h+q^{-A}\right)\right).
\end{equation}
\item In either case the following finer additive estimate holds (with absolute $c>0$):
\begin{equation}\label{eq:additive}
 \normcount=1+s_k\chi_*(g)b^t
   +O\left(\frac{\lambda^2 e^{-c\lambda}}{h}+q^{-A-2}\right).
\end{equation}
The implied constants in the finite-sampling error can be absolute; the onset of the arithmetic estimates depends on $a,A,C,B$.
\end{enumerate}
\end{theorem}

In contrast to a result requiring $m/N>1/2+\delta$, this result allows any fixed density above $1/3$ and all distances $t$, even one exceptional prime. It quantifies what can replace a density margin: one needs enough minority signs in a typical sample, not just a positive distance from a quadratic fibre.

\begin{corollary}[Exact equidistribution criterion]\label{cor:criterion}
In the setting of Theorem~\ref{thm:main}, along any sequence $q\to\infty$, define
\[
 D_{S,k}=\max_{g\in G}\abs{\frac{\varphi(q)R_{S,k}(g)}{\binom mk}-1}.
\]
Then $D_{S,k}=b^t+O(h^{-1}+q^{-A})$. Consequently $D_{S,k}\to0$ if and only if $b^t\to0$.
If $h=o(m)$, this is equivalent to $ht/m\to\infty$. In that sparse regime:
\begin{align*}
 ht/m\to\lambda_0\in(0,\infty)&\implies D_{S,k}\to e^{-2\lambda_0};\\
 t\ge1,\ ht/m\to0&\implies
 \frac{\varphi(q)R_{S,k}(g)}{\binom mk}\sim 2ht/m
 \quad\text{when }s_k\chi_*(g)=-1.
\end{align*}
If $h/m\to v\in(0,1/2]$ and $t$ is fixed and positive, the limit is $(1-2v)^t$; thus nearly balanced quotas can mix even with only one minority prime. The total variation distance of the product law from uniform is $\tfrac12 b^t+O(h^{-1}+q^{-A})$.
\end{corollary}

\section{An unconditional finite-sampling theorem}
The key is a parity estimate with no growing prefactor. For integers $m\ge2$ and $0\le h,t\le m/2$, let $J$ count the marked elements in a uniform $h$-subset of an $m$-set with $t$ marked elements. Define
\begin{equation}\label{eq:parity}
 F(m,h,t)=\E(-1)^J=\frac{\sum_j(-1)^j\binom tj\binom{m-t}{h-j}}{\binom mh}.
\end{equation}
All out-of-range binomial coefficients are zero. Symmetry of the intersection of two subsets gives $F(m,h,t)=F(m,t,h)$, so $J$ may also be viewed as the number of successes in $t$ draws from a population with $h$ successes.

\begin{lemma}[Uniform exponential parity decay]\label{lem:decay}
There is an absolute constant $K$ such that for $1\le h\le m/2$, $0\le t\le m/2$,
\[
 |F(m,h,t)|\le K\exp\left(-\frac{ht}{4m}\right).
\]
\end{lemma}
\begin{proof}
For $m<8$, enlarge $K$; assume $m\ge8$. Pair $r=\min(t,\floor{m/4})$ marked elements with distinct unmarked elements. This is possible since $t\le m/2$. Say a pair splits if exactly one member is selected. Conditional on the selection outside the pairs and on the number selected in each pair, the orientations of split pairs are uniform and independent, by the symmetry of uniform sampling. If any pair splits, the conditional mean of $(-1)^J$ is zero. Hence $|F|\le\Pp(\text{no split})$.

To bound that probability, select each element independently with probability $u=h/m$, and then condition on total count $B_m=h$. Write $E$ for no split. Under independent sampling, $\Pp(E)=(1-2u(1-u))^r\le e^{-ur}$. Conditional on $E$, the unpaired $m-2r$ elements are still independent Bernoulli variables, also independent of the counts in the pairs. Therefore
\[
 \Pp(B_m=h\mid E)\le \sup_j\Pp(\operatorname{Bin}(m-2r,u)=j)\ll h^{-1/2}.
\]
For completeness, for $n=m-2r\ge m/2$ the binomial atom bound follows by Fourier inversion:
\[
 \sup_j\Pp(\operatorname{Bin}(n,u)=j)
 \le \frac1{2\pi}\int_{-\pi}^{\pi}e^{-2nu(1-u)\sin^2(\theta/2)}\,d\theta
 \ll [nu(1-u)]^{-1/2}\ll h^{-1/2}.
\]
Here $1-u\ge1/2$ and $|\sin(\theta/2)|\ge|\theta|/\pi$ on the integration interval. The binomial law $\operatorname{Bin}(m,u)$ has a mode at its integer mean $h$ (successive atom ratios prove this); its variance is at most $h$. Chebyshev's inequality places at least $3/4$ of its mass in an interval of length $4\sqrt h$ about $h$, whence $\Pp(B_m=h)\gg h^{-1/2}$. Bayes' formula now gives $\Pp(E\mid B_m=h)\ll e^{-ur}$. Since $r\ge t/4$, the claim follows.
\end{proof}

\begin{theorem}[Localized binomial approximation; unconditional]\label{thm:parity}
For $2\le h\le m/2$, $0\le t\le m/2$, $u=h/m$, $b=1-2u$ and $\lambda=ut$,
\begin{equation}\label{eq:localpar}
 \abs{F(m,h,t)-b^t}\ll \frac{\lambda^2 e^{-\lambda/16}}{h}
 \ll \frac{\min(1,\lambda)}{h}.
\end{equation}
All constants are absolute. At $t=0,1$ the difference vanishes exactly. In the limit $h\to\infty$, $h/m\to0$, $ht/m\to\lambda_0\in(0,\infty)$,
\begin{equation}\label{eq:sharp}
 h\bigl(F(m,h,t)-b^t\bigr)\longrightarrow -2\lambda_0^2e^{-2\lambda_0}.
\end{equation}
\end{theorem}
\begin{proof}[Proof of the localized approximation]
Let $J_i$ be the number of successes in the first $i$ draws without replacement from the $m$-population with $h$ successes, and $F_i=\E(-1)^{J_i}=F(m,h,i)$. The next draw has conditional success probability $(h-J_i)/(m-i)$, so
\begin{equation}\label{eq:recurrence}
 F_{i+1}=bF_i+\frac{2}{m-i}\E\bigl[(-1)^{J_i}(J_i-iu)\bigr].
\end{equation}
Exchangeability, followed by conditioning on a specified successful draw, gives the exact identity
\[
 \E[J_i(-1)^{J_i}]=-iu\,F(m-1,h-1,i-1)\quad(i\ge1).
\]
By Lemma~\ref{lem:decay}, for $i\ge2$ and $h\ge2$,
\[
 |F_i|\ll e^{-ui/4},\quad
 |F(m-1,h-1,i-1)|\ll e^{-ui/16}.
\]
Indeed $(h-1)(i-1)/(m-1)\ge hi/(4m)$; the smaller population still satisfies both half-size restrictions. For $i=1$ the same estimate holds after enlarging the constant because the shifted bias is one and $u\le1/2$. The $i=0$ error is zero. Since $i<t\le m/2$, the correction in \eqref{eq:recurrence} is at most $K_1 (ui/m)e^{-ui/16}$ in modulus. Iterating from $F_0=1$ and using $0\le b\le e^{-2u}$ yields
\begin{align*}
 |F_t-b^t|&\le \frac{K_1u}{m}\sum_{i=1}^{t-1}i\,e^{-ui/16}b^{t-1-i}\\
 &\le \frac{K_1u}{m}e^{-u(t-1)/16}\sum_{i=1}^{t-1}i
 \ll \frac{\lambda^2}{h}e^{-\lambda/16}.
\end{align*}
The last identity uses $ut^2/m=\lambda^2/h$. For $0\le\lambda\le1$, $\lambda^2\le\lambda$, while for $\lambda\ge1$, $\lambda^2e^{-\lambda/16}$ is bounded. This proves \eqref{eq:localpar}.
\end{proof}
\begin{proof}[Proof of the sharp first-order correction]
Write $(n)_j=n(n-1)\cdots(n-j+1)$. Expand $(-1)^{J_t}=(1-2)^{J_t}$ and take expectations. By symmetry the factorial moments satisfy
\[
 F(m,h,t)=\sum_{j=0}^{t}\binom tj(-2)^j\frac{(h)_j}{(m)_j},\qquad
 b^t=\sum_{j=0}^{t}\binom tj(-2u)^j.
\]
For $j\le t$, set $r_j=(h)_j/((m)_j u^j)$; put $r_j=0$ if $j>h$. For $j\le h$,
\[
 r_j=\prod_{\ell=0}^{j-1}\frac{1-\ell/h}{1-\ell/m},\qquad
 0\le1-r_j\le\sum_{\ell<j}\frac{\ell/h-\ell/m}{1-\ell/m}\le\frac{j(j-1)}h.
\]
For $j>h$ the final bound still holds. For fixed $j$, expansion of the finite product gives
$h(r_j-1)\to-j(j-1)/2$, as $h/m\to0$. Since $ut\to\lambda_0>0$, we have $t\to\infty$, $\binom tj u^j\to\lambda_0^j/j!$, and for all $j\le t$ we have $\binom tj u^j\le(ut)^j/j!$. Thus the summands of $h(F-b^t)$ are dominated by a summable sequence $j(j-1)(2M)^j/j!$ for some fixed $M>\sup ut$ (after discarding finitely many terms of the sequence of populations). Dominated convergence gives
\[
 \lim h(F-b^t)=-\tfrac12\sum_{j\ge0}j(j-1)\frac{(-2\lambda_0)^j}{j!}
 =-2\lambda_0^2e^{-2\lambda_0}.
\]
In particular the uniform order $1/h$ cannot be improved to $o(1/h)$, even restricting to the sparse critical window.
\end{proof}

\section{Arithmetic input and isolation of the exceptional character}
\begin{lemma}[Small bias at the endpoint; uses DGRH]\label{lem:bias}
Given $\epsilon>0$, there is $C>0$ such that, at $y=C\log^2 q$ and for all sufficiently large $q$,
\[
 N\sim\frac{y}{\log y},\qquad
 \max_{\chi\ne1}\abs{\frac1N\sum_{p\in\Pool}\chi(p)}\le\epsilon.
\]
\end{lemma}
\begin{proof}
We give a smoothed explicit-formula proof of the article's lemma \texttt{dir:lem:bias}. Fix $0<\delta<1/4$ and choose $W\in C_c^\infty((0,1))$ with $0\le W\le1$, equal to one on $[\delta,1-\delta]$. Its Mellin transform $\widetilde W(s)=\int W(v)v^{s-1}\,dv$ decays faster than every inverse power of $|\Im s|$ on fixed vertical strips. For a primitive nonprincipal $\psi$ of conductor $f\le q$, the smoothed explicit formula gives
\begin{equation}\label{eq:smoothed}
 \sum_{n\ge1}\Lambda(n)\psi(n)W(n/y)
 =-\sum_{\rho}y^\rho\widetilde W(\rho)+O_W(y^{1/4}\log(2f)).
\end{equation}
Here $\rho$ runs over the nontrivial zeros, with multiplicity. We recall the standard analytic inputs \cite[Chs.~10, 12--13]{MV}: Mellin inversion of $-L'/L$, the local logarithmic derivative estimate away from zeros, and the bound $O(\log(2f(2+|T|)))$ for the number of zeros in a unit-height strip. More explicitly, start at $\Rea s=2$ and shift to $1/4$ using \eqref{eq:dgrh}; the functional equation and local logarithmic derivative estimate bound $L'/L(1/4+iv,\psi)$ by $O(\log(2f(2+|v|)))$ because all nontrivial zeros have real part $1/2$. The shift crosses no pole of $L$ or trivial zeros, and horizontal integrals vanish on heights avoiding zero ordinates by inverse-logarithmic distances, using the same zero-count bound and rapid decay. This proves \eqref{eq:smoothed}. The zero sum is $O_W(\sqrt y\log(2f))$ by \eqref{eq:dgrh} and the unit-interval zero-count estimate. Removing prime powers costs $O(\sqrt y\log^2 y)$ (sum over integer bases and exponents), and replacing $\psi$ by its induced character modulo $q$ in the prime sum costs at most $\sum_{p\mid q}\log p\le\log q$. Hence, uniformly for all nonprincipal $\chi$ modulo $q$,
\[
 \abs{\sum_{p\in\Pool}\chi(p)(\log p) W(p/y)}
 \ll_W \sqrt y\log(2q)+\sqrt y\log^2 y+\log q.
\]
By the prime number theorem, replacing $W$ by one costs at most $(2\delta+o(1))y$ in absolute value. Replacing $\log p$ by $\log y$ costs at most $\pi(y)\log y-\vartheta(y)=o(y)$. Finally at most $\log q/\log 2$ primes divide $q$, so $N\log y\sim y$. On division by $N\log y$ the bound is at most $2\delta+O_W(C^{-1/2})+o(1)$. First choose $\delta$ small and then $C$ large. The constants are uniform in $q,\chi$; the onset may depend on the fixed choices.
\end{proof}

\begin{lemma}[Only one possible large nonprincipal mode]\label{lem:isolation}
Fix $a>1/3$. There are $\tau,\epsilon>0$ depending only on $a$ such that the following holds for every finite pool of labelled elements in a finite abelian group with all nonprincipal character means at most $\epsilon$ in modulus. If a set $S$ retains a fraction at least $a$ of the pool, at most one nonprincipal character has bias $>1-\tau$ on $S$, and that character, if present, is quadratic.
\end{lemma}
\begin{proof}
With the normalized counting inner product on the full pool, distinct character vectors have pairwise inner product at most $\epsilon$ in modulus. Thus any three distinct character vectors have Gram matrix of operator norm at most $1+2\epsilon$, by the row-sum bound. Put $\alpha=m/N$. If each of the three has bias $>1-\tau$ on $S$, then their squared inner products with the indicator of $S$ sum to more than $3\alpha^2(1-\tau)^2$, yet this sum is at most $(1+2\epsilon)\alpha$ by the Gram operator bound. Choose $\tau,\epsilon$ with $3a(1-\tau)^2>1+2\epsilon$. The principal character has bias one; hence there cannot be two distinct nonprincipal characters of bias $>1-\tau$. A nonreal character and its inverse have equal bias, and if a nonprincipal character equals its inverse then its order is two. The claim follows.
\end{proof}

\begin{lemma}[Fixed quota Fourier estimate]\label{lem:quota}
For any complex numbers $z_1,\dots,z_m$ of modulus one, $1\le k<m$, put $v=|\sum z_j|/m$. Then
\[
 \abs{\binom mk^{-1}[Z^k]\prod_{j=1}^m(1+Zz_j)}
 \le(m+1)\exp\left(-(1-v)\frac{k(m-k)}m\right).
\]
\end{lemma}
\begin{proof}
On the circle $|Z|=r=k/(m-k)$, use
\[
 |1+re^{i\theta}z_j|^2=(1+r)^2\left(1-\frac{2r(1-\Rea(e^{i\theta}z_j))}{(1+r)^2}\right).
\]
The inequality $1-x\le e^{-x}$ gives the product bound $(1+r)^m\exp(-rm(1-v)/(1+r)^2)$. Cauchy's coefficient bound divides by $r^k\binom mk$. The remaining factor $(1+r)^m/(r^k\binom mk)$ is the reciprocal of the probability at the integer mean of $\operatorname{Bin}(m,k/m)$, at most $m+1$ since the mean is a mode. This is the article's circle argument \cite[\texttt{dir:lem:circle}]{Article}.
\end{proof}

\section{Proof of the arithmetic theorem and sharpness}
\begin{proof}[Proof of Theorem~\ref{thm:main}]
Choose $\tau,\epsilon$ by Lemma~\ref{lem:isolation}, then $C$ by Lemma~\ref{lem:bias}. Since $m\le C L^2$ and $h\ge BL$, choosing $B$ sufficiently large in terms of $\tau,A$ gives
\begin{equation}\label{eq:remainder}
 q(m+1)e^{-\tau h/2}\le q^{-A-2}.
\end{equation}
For example $B=2(A+5)/\tau$ suffices for all sufficiently large $q$. For $\chi\in\widehat G$ let $E_k(\chi)$ be the mean of $\chi(\prod_{p\in T}p)$ over uniform $k$-subsets $T$. By minimizing $t$ in \eqref{eq:distance}, $\chi_*$ maximizes quadratic bias (the maximum is $1-2t/m$), so Lemma~\ref{lem:isolation} shows every nonprincipal character except $\chi_*$ has bias at most $1-\tau$. Fourier inversion and Lemma~\ref{lem:quota}, with $k(m-k)/m\ge h/2$, give the exact one-mode reduction
\begin{equation}\label{eq:onemode}
 \normcount=1+\chi_*(g) E_k(\chi_*)+\mathcal E_g,
 \qquad |\mathcal E_g|\le q^{-A-2}.
\end{equation}
This holds regardless of ties among minimizing characters. The quadratic mean is exactly $s_k F(m,h,t)$: for $k\le m/2$, each factor contributes its majority sign times the minority parity; for $k>m/2$, the selected minority count equals $t$ minus the omitted minority count. Thus \eqref{eq:additive} follows from Theorem~\ref{thm:parity}.

If $t=0$, every product has sign $\sigma^k=s_k$. The other sign class is exactly empty and \eqref{eq:onemode} proves the stated count in the occupied class. If $t\ge1$, the smaller model value obeys
\begin{equation}\label{eq:minprofile}
 1-b^t\ge1-e^{-2\lambda}\ge (1-e^{-2})\min(1,\lambda).
\end{equation}
The last inequality follows from concavity on $[0,1]$ and monotonicity afterwards. By \eqref{eq:localpar} the finite-sampling error relative to either profile is $O(1/h)$. Since $\lambda\ge h/m$, dividing the residual in \eqref{eq:onemode} by the smaller profile costs at most $O(m/h)=O(L^2)$, giving $O(q^{-A})$. Both relative errors tend uniformly to zero. Since both profiles are strictly positive, every class occurs for sufficiently large $q$. This proves \eqref{eq:relative} and the theorem.
\end{proof}

\begin{proof}[Proof of Corollary~\ref{cor:criterion}]
Both signs occur in $G$, and their classes have equal cardinality. The uniform additive estimate and the reverse triangle inequality show $D_{S,k}=b^t+O(1/h+q^{-A})$ and, by averaging absolute deviations, the asserted total variation formula. If $u\to0$, then $\log(1-2u)=-2u+O(u^2)$, so $b^t\to0$ exactly when $ut\to\infty$ (including sequences with $t=0$). If $ut\to\lambda_0\in(0,\infty)$, $b^t\to e^{-2\lambda_0}$. For $t\ge1$ and $ut\to0$, $1-b^t\sim2ut$; apply the relative estimate, which remains valid even on this sparse sign class. The fixed-$t$ dense case follows by continuity.
\end{proof}

\begin{corollary}[Sharp correction in the arithmetic critical window]\label{cor:correction}
In Theorem~\ref{thm:main}, suppose $h=o(m)$ and $ht/m\to\lambda_0\in(0,\infty)$. Then, uniformly in $g$,
\[
 \normcount=1+s_k\chi_*(g)\left(b^t-\frac{2\lambda_0^2e^{-2\lambda_0}}h\right)+o(1/h).
\]
In particular $h(D_{S,k}-b^t)\to-2\lambda_0^2e^{-2\lambda_0}$.
\end{corollary}
\begin{proof}
Apply \eqref{eq:sharp} in \eqref{eq:onemode}; the residual is $o(1/h)$ since $h\le m=O(L^2)$. The limiting bias is positive, so taking the maximum over both signs gives the second conclusion.
\end{proof}

\begin{proposition}[Realization of the transition and sharpness of the density floor]\label{prop:sharp}
At the logarithm-squared cutoff with sufficiently large fixed $C$, all finite transition values $e^{-2\lambda_0}$, $\lambda_0\ge0$, are realized by retained sets of density at least $1/2$ with $h=\ceil{B\log q}$, along prime moduli. Also, at retention density $1/3$, a criterion based only on distance from quadratic fibres fails, even when that distance is a fixed positive fraction of $m$.
\end{proposition}
\begin{proof}
For prime $q$, there is just one nonprincipal quadratic character. Let $S_0$ be its larger sign class in $\Pool$, of size $m_0\ge N/2$. The other class has size at least $(1-\epsilon)N/2$ by Lemma~\ref{lem:bias}. Set $h=k=\ceil{BL}$ and add any $t=\floor{\lambda_0 m_0/h}$ primes from the other class. Since $N\sim C L^2/\log(C L^2)$, $h=o(m_0)$ and $t=o(m_0)$; this choice is feasible for all sufficiently large $q$. Then $ht/(m_0+t)\to\lambda_0$. For $\lambda_0=0$ one may choose either no minority primes or just one, the latter exhibiting full support but limiting deviation one. Adding instead $t=\floor{(m_0/h)\sqrt h}$ yields $ht/m\to\infty$ and limit zero. This proves sharpness of the critical scale even when the retained density is at least $1/2$.

For the floor, choose arbitrarily large prime $q\equiv1\pmod6$ and a character $\eta$ of order three. Choose $S$ of size $m=\ceil{N/3}$ inside its largest fibre in $\Pool$. The unique quadratic character $\chi$ and $\eta$ generate six distinct characters. Orthogonality on the quotient of order six, together with small bias $\epsilon<1/10$ for the five nontrivial characters, shows that each joint fibre has size at most $(1+5\epsilon)N/6$. Therefore either sign class of $\chi$ in $S$ has size at least $(1-5\epsilon)N/6$, and $t/m\ge(1-5\epsilon)/6>0$. Nevertheless all $k$-products lie in the single fibre $\eta(g)=\eta(S)^k$, omitting two thirds of the units. The $1/3$ floor is thus sharp for a purely quadratic obstruction theorem. This does not rule out a higher-order quotient law at or below the floor.
\end{proof}

\section{Interpretation, novelty, and scope}
The source article's fixed positive majority margin forces a constant Fourier gap and a power-saving error in all classes. Here the exceptional quadratic coefficient is retained, not discarded. The new unconditional estimate \eqref{eq:localpar}, with an error proportional to the smaller profile, is what upgrades the existing exact coefficient to a uniform relative asymptotic, including a sign class whose mass tends to zero. The sharp correction \eqref{eq:sharp} also distinguishes sampling without replacement from independent signs. At sparse quotas $h\asymp\log q$, one has $m\asymp (\log q)^2/\log\log q$ uniformly at positive density, so the crossover minority count is of order $\log q/\log\log q$. This is far smaller than a fixed fraction of the pool. In particular one exceptional prime guarantees support (already implied by the article), but not equidistribution; a growing number on the crossover scale still leaves a nonzero bias.

One must distinguish asymptotic equidistribution from a prescribed power saving. Equation~\eqref{eq:onemode} is power-saving in every nonexceptional mode. Lemma~\ref{lem:decay} shows that the exceptional mode is also $O(q^{-A})$ if $ht/m\ge4(A+1)\log q$. A fixed positive minority fraction therefore recovers a power saving with a large enough quota. The condition $ht/m\to\infty$ alone yields $o(1)$, not in general $O(q^{-A})$. For instance, using the construction of Proposition~\ref{prop:sharp} with $t=\floor{(m_0/h)\log\log h}$, at sparse quotas and $ut=\log\log h+o(1)$, $b^t\asymp (\log h)^{-2}$ along the logarithm-squared regime with $h\asymp L$, while the correction in \eqref{eq:localpar} is $O((\log\log h)^2/h)$, negligible. Thus no fixed power of $q$ can bound the deviation in this example. The article's question is answered in the qualitative and quantitative-profile senses, not by a claim of power saving for arbitrary slowly diverging distance.

The related work of Booker--Pomerance \cite{BP} and Munsch--Shparlinski \cite{MS} concerns small smooth squarefree representatives. Jao--Miller--Venkatesan \cite{JMV} use GRH character estimates for small-prime Cayley graph expansion, where walks may repeat generators. Klurman--Shparlinski--Ter\"av\"ainen \cite{KST} study subset products of small primes for almost all moduli, and Pach \cite{Pach} studies fixed-cardinality subset sums in a finite abelian group. The merged ProveIt article already establishes fixed-quota support and relative counts via bounded quotients under DGRH. Our more specific contribution is the uniform binomial crossover with exponentially localized $1/h$ error, the first-order correction, and the necessary-and-sufficient distance criterion in the sparse regime. A targeted comparison with those works and the relevant article sections did not identify this exact combined result; an exhaustive literature or all-914-result novelty certification was not performed. We claim a new derivation and a candidate research contribution, not verified global priority.

\paragraph{Limitations.} The constants $C,B$ are existential, not optimized numerical constants. No theorem is claimed for $h=o(\log q)$, which is another open question in the article, or at density $\le1/3$ using only quadratic information. We have not formalized the finite probability lemmas or the uniform explicit formula in Lean, and did not rerun the vendored build. The supplied exact integer computations are checks of finite identities and illustrative counts, not proof of any asymptotic assertion.

\section{Reproducible checks}
The companion \texttt{verify.py} uses only the Python standard library. It enumerates all triples $2\le m\le80$, $1\le h\le m/2$, $0\le t\le m/2$ (44,280 triples), checking the signed coefficient formula, symmetry, complementary quotas, recurrence \eqref{eq:recurrence}, and boundary cases exactly over rational numbers. It also evaluates the first-order correction at $m=h^2$, $t=\lambda_0 h$, and computes residue counts by exact integer dynamic programming for $q=101$, $y=251$ across sign imbalances. The latter checks aggregate quadratic profiles exactly; at small moduli the nonquadratic residual need not be small.

\begin{center}
\begin{tabular}{rrrr}
\toprule
$\lambda_0$ & $h$ & $h(F-b^t)$ & limiting correction\\
\midrule
0.5 & 640 & -0.1840597 & -0.1839397 \\
1.0 & 640 & -0.2708120 & -0.2706706 \\
2.0 & 640 & -0.1458746 & -0.1465251 \\
\bottomrule
\end{tabular}
\end{center}
The data tables and run transcript are included in the archive. The report and scripts do not modify the ProveIt repository.

\begin{thebibliography}{99}
\bibitem{Article} ProveIt contributors, \emph{Beyond the Riemann hypothesis for Dirichlet $L$-functions: consequences, transfers and limits}, 26 Nov. 2026. Sections \texttt{src.tex} and \texttt{dir.tex}, especially labels \texttt{dir:lem:bias}, \texttt{dir:lem:circle}, \texttt{dir:thm:quotient}, \texttt{dir:sec:open}. See \href{https://github.com/VladimirReshetnikov/ProveIt/tree/main/Analysis/GeneralizedRiemannHypothesis/Article}{article repository}.
\bibitem{Lean} L. Sernau and collaborators, \emph{The Generalized Riemann Hypothesis for Dirichlet L-functions, in Lean 4}, upstream revision \texttt{72fead102f663d5c7a407964a13dc0cf9d977207}, 25 Nov. 2026. \href{https://github.com/lsernau/grh-lean/tree/72fead102f663d5c7a407964a13dc0cf9d977207/grh}{Pinned source}; \href{https://github.com/VladimirReshetnikov/ProveIt/tree/main/Analysis/GeneralizedRiemannHypothesis/vendor/grh-lean}{vendored copy}.
\bibitem{Suth} A. Sutherland, \emph{Verification files for the shorter Lean proof of RH and of GRH for Dirichlet L-functions}, updated 26 Nov. 2026, \url{https://klein.mit.edu/~drew/shortgrh/}.
\bibitem{MV} H. L. Montgomery and R. C. Vaughan, \emph{Multiplicative Number Theory I: Classical Theory}, Cambridge University Press, 2007.
\bibitem{BP} A. R. Booker and C. Pomerance, Squarefree smooth numbers and Euclidean prime generators, \emph{Proc. Amer. Math. Soc.} \textbf{145} (2017), 5035--5042. \href{https://doi.org/10.1090/proc/13576}{doi:10.1090/proc/13576}.
\bibitem{MS} M. Munsch and I. E. Shparlinski, On smooth square-free numbers in arithmetic progressions, \emph{J. London Math. Soc.} \textbf{101} (2020), 1041--1067. \href{https://doi.org/10.1112/jlms.12297}{doi:10.1112/jlms.12297}.
\bibitem{JMV} D. Jao, S. D. Miller, R. Venkatesan, Expander graphs based on GRH with an application to elliptic curve cryptography, \emph{J. Number Theory} \textbf{129} (2009), 1491--1504. \href{https://arxiv.org/abs/0811.0647}{arXiv:0811.0647}.
\bibitem{KST} O. Klurman, I. E. Shparlinski, J. Ter\"av\"ainen, On the generation of multiplicative groups by small primes, preprint, 2026. \href{https://arxiv.org/abs/2609.25393}{arXiv:2609.25393}.
\bibitem{Pach} P. P. Pach, Equidistribution of subset sums, preprint, 2025. \href{https://arxiv.org/abs/2505.12496}{arXiv:2505.12496}.
\end{thebibliography}
\end{document}
